% !TeX root = ../papers/nc-p3.tex
\section{四つの異なる文字に支えられた窓の全分類と二色欠損への適用}
\label{sec:squarefree-window-coverage}
\begin{theorem}\label{thm:squarefree-window-coverage}
四位置に四ラベルの基底があり、最上テンソルの台が相異なる四文字の置換に含まれるとする。
隣接階数がともに6なら、標準Koszul左右AS正則であるための必要十分条件は
二十四の置換係数が全て非零であることである。零係数がある全ての場合で指数増大する。
正則窓は添字保持同型まで量子多項式窓の全体に一致する。最小周期は1であり、
相互の同型は系\ref{cor:quantum-torus-normal-form}の三変数と $S_4$ で決まる。

特にトーラスが四位置に作用してテンソル直線を保ち、各位置の四つの重みが一斉平行移動を除き
同じアフィン独立な四点になる全ての標準窓は、この量子枝に属する。
孤立単純印や反対切断階数には仮定を課さない。
\end{theorem}
\begin{proof}
各位置の隣接切断は第一対の無順序二文字を固定した六つの二行二列ブロックに分かれる。
一つでも零なら $ij,ji$ は共にその位置の関係の台に現れない。
全位置で $ii,jj$ も関係に現れないため、対角で待つ語を挿入して
$i\to j\to i$ の閉路と一定の閉路を同じ長さで反復でき、指数増大する単項式商を得る。
従って非指数の場合は六ブロックとも非零、切断階数6から各階数1である。

各二行二列の非零階数1行列の台は九通りの非空長方形のいずれかである。
最初の切断の六台を全て選び、次の切断の六台も非空長方形であるものを残すと、
$9^6=531441$ 通りから22785通りとなる。これらについて、ある位置で
$ij$ を行とする二成分がともに零のときに辺 $i\to j$ を引く。
閉路があれば対角語による待機を使った同じ単項式商が指数増大する。
四頂点の推移閉包で対角が一つもないものは二十四成分全て非零の一台だけである。
\path{computations/verify_squarefree_window_coverage.py} は四文字の置換から
十二ブロックの位置を生成し、全ての真理値をこの手順で検査する。

残る台では全位置 $p$ で各無順序対に一つだけ
$x_{p,j}x_{p+1,i}=q_{p,ij}x_{p,i}x_{p+1,j}$ ($i<j$) があり $q_{p,ij}\ne0$。
最上テンソルの係数で隣接二文字を交換する比は同じ関係直線で決まり、残り二文字の順序に依存しない。
三文字の反転を二通りに行い、全係数非零を使って比較すると
\[
   q_{p,ij} q_{p+1,ik} q_{p,jk}
       =q_{p+1,ij}q_{p,ik}q_{p+1,jk}\qquad(i<j<k)
\]
である（関係直線と書換係数の間の負号は両辺で同数）。
よって各 $p$ の $q_{p,ij}/q_{0,ij}$ は四ラベルの倍率の比に書ける。
位置 $p$ の各文字に掛ける倍率と次位置の倍率との比を順に指定し、両方向に延ばせば
全係数を一定の $q_{0,ij}$ にできる。これは添字を動かさない基底変更である。
逆に全ての量子多項式窓は非零の歪交代テンソルをこの台に持つ。
トーラスの場合は重みの平行移動分を除き、アフィン独立性によりテンソルの全項で
各文字の出現回数が一定である。各位置の一次縮約は単射なので各文字が少なくとも一度現れる。
四文字で長さ4のため全て一度ずつであり、本定理に帰着する。
\end{proof}

\begin{proposition}\label{prop:klein-unmarked-matching-reduction}
四元群次数零の四テンソルで両隣接切断の色 $0$ のブロックが零、他の三色が各階数2とする。
色の原点は四位置のラベルの平行移動で選んでよい。四文字が相異なる全二十四成分が非零とする。
相異なる二ラベル $i,j$ について $d_{ij}=w_{ijij}$ と置くと、$d_{ij}=0$ と $d_{ji}=0$ は同値である。
四頂点の無向グラフに $d_{ij}\ne0$ の対を辺として取る。
\begin{enumerate}
\item 次数3の頂点があれば、四元群の印付き交換子表示があり、前三節の全点判定が適用される。
\item 互いに素な二辺がなければ必ず標準Koszul左右AS正則で、量子窓を軌道閉包に持つ。
\item 残るグラフは完全マッチング、四頂点の道、四辺のサイクルの三種類だけである。
道の中央の辺は全置換係数を変えずに対角基底変更の極限で消すことができる。
\end{enumerate}
\end{proposition}
\begin{proof}
二つの欠色と総次数により、置換以外で許されるのはちょうど $ijij$ である。
$\{a,b\},\{c,d\}$ を一つの完全マッチングとして色 $a+b$ のブロックをこの二組の順に並べると、
対角側は $\operatorname{diag}(d_{ab},d_{ba})$ と $\operatorname{diag}(d_{cd},d_{dc})$、
非対角側は全成分非零の二行二列行列 $U,V$ となる。
二つの対角位置 $a,b$ と、残る組で行と列を異ならせた三次小行列式は
$-d_{ab} U_{b d} V_{c b}-d_{ba} U_{a d} V_{c a}$ である。
その消失と両係数の非消失で最初の同値が従う。
次数3の頂点を星とすれば対向星成分の二つの三次対角行列が可逆である。それらを単位に揃えた
Schur補元が印付き交換子表示であり、四元群次数から各行列は残る二軸の間の二成分に限る。

全四位置で同じラベル重み $v_i$ を用い、テンソルを総重み $\sum_i v_i$ で割れば
置換成分は固定される。対 $ab$ の相対重みは $v_a+v_b-v_c-v_d$ である。
三つのマッチングの相対重みは独立で、補対の相対重みはその負である。
従って互いに素な非零辺がなければ全非零辺に正の整数重みを与えられる
（有理解の分母を払う）。極限には全置換成分が残り、各隣接切断には六つの非零ブロックがあり、
元の階数6以下なので全て階数1である。定理\ref{thm:squarefree-window-coverage}と
命題\ref{prop:unrestricted-good-reduction}より元の窓も全次数で正則となる。
最大次数2でマッチングを含む四頂点グラフの三種類は直接列挙できる。
道では二つの端辺の重みを零、中央の辺の重みを正とすればよい。
\end{proof}
